\documentclass[11pt]{article}

\usepackage[T1]{fontenc}
\usepackage[utf8]{inputenc}
\usepackage{amsmath,amssymb,amsthm}
\usepackage[a4paper,margin=30mm]{geometry}
\usepackage[hidelinks,
  pdftitle={The Two-Hole Minimum in the Regular Hyperbolic (4,5)-Tessellation Is 17},
  pdfauthor={Xiaojian Zeng}]{hyperref}
\usepackage{booktabs}

\newtheorem{theorem}{Theorem}
\newtheorem{lemma}[theorem]{Lemma}

\title{The Two-Hole Minimum in the Regular Hyperbolic
  $\{4,5\}$-Tessellation Is 17}
\author{Xiaojian Zeng\\
  \small PBC School of Finance, Tsinghua University, Beijing 100190, China\\
  \small \texttt{xiaojianzeng@foxmail.com}}
\date{10 September 2026}

\begin{document}

\maketitle

\begin{abstract}
Let $g_{4,5}(h)$ be the minimum number of square tiles in an
edge-connected polyform with exactly $h$ holes in the regular hyperbolic
$\{4,5\}$-tessellation.  Eldridge et al. gave a $17$-tile example with two
holes and a perimeter obstruction below $16$, leaving the endpoint $16$
undecided.  We prove that no $16$-tile example exists, and hence that
$g_{4,5}(2)=17$.  Equality in their perimeter bound would force the two
holes to be single tiles and their filling to be an $18$-tile
minimum-perimeter animal.  A uniqueness theorem of Malen, Rold\'an, and
Toal\'a-Enr\'iquez identifies that animal with the second layer grown from a
domino.  Its only fully edge-surrounded tiles are the adjacent central pair,
whose simultaneous deletion makes one cavity rather than two.  An exact
Coxeter-group certificate over $\mathbb Q(\sqrt2,\sqrt5)$ verifies the finite
incidence facts used in this last step.
\end{abstract}

\noindent\textbf{Keywords.}
Hyperbolic polyforms; holes; extremal animals; isoperimetric inequality;
exact certificate.

\section{The endpoint problem}

Let $\mathcal T$ be the tile set of the regular $\{4,5\}$-tessellation of
$\mathbb H^2$.  Thus every tile is a regular geodesic square and five tiles meet
at every tessellation vertex.  A \emph{polyform} is a finite set
$A\subset\mathcal T$ whose edge-dual graph is connected: two tiles are
adjacent precisely when they share a complete tessellation edge.  Tiles may
also have vertex-only contacts.  A hole is a bounded connected component of
$\mathbb H^2\setminus\bigcup_{T\in A}T$, and
\[
 g_{4,5}(h)=\min\{\lvert A\rvert:A\text{ is a polyform with exactly }h
 \text{ holes}\}.
\]
These are the conventions of Eldridge et al.~\cite{EldridgeEtAl2026}.

Their Theorem~1.2 and Table~1 supply a $17$-tile polyform with two holes.
For this entry, their Table~2 gives
\[
 M_{4,5}(15,2)=1.5,
 \qquad M_{4,5}(16,2)=2,
 \qquad M_{4,5}(17,2)=2.
\]
Lemma~2.1 and Corollary~2.2 of the same paper therefore rule out $15$ tiles
but not $16$, giving the documented window
\[
                 16\le g_{4,5}(2)\le17.
\]
We settle precisely this one-unit alternative.

\begin{theorem}\label{thm:main}
There is no $16$-tile $\{4,5\}$-polyform with exactly two holes.
Consequently,
\[
                         g_{4,5}(2)=17.
\]
\end{theorem}

The proof combines the equality case of the preceding perimeter bound with
the unique-extremal subsequence of Malen, Rold\'an, and
Toal\'a-Enr\'iquez~\cite{MalenRoldanToala2023}.  It does not classify
$\{4,5\}$-animals at arbitrary sizes and makes no general assertion about
$g_{p,q}(h)$.

\section{What equality at 16 would force}

For a polyform $A$, write $b(A)$ for the number of shared tile edges,
$p_o(A)$ for its outer perimeter, and $p_h(A)$ for the total perimeter of
its holes.  Let $\overline A$ be obtained by filling every hole with all the
tiles it contains.  Finally, let $\mathcal P_{\min}(N)$ denote the minimum perimeter of
an $N$-tile $\{4,5\}$-animal.

The function used in~\cite{EldridgeEtAl2026} is
\[
 M_{p,q}(n,h)
   =\frac{n(p-2)+2-\mathcal P_{\min}^{p,q}(n+h)}{p}.
\]
Substituting $p=4$, $n=16$, $h=2$, and the tabulated value
$M_{4,5}(16,2)=2$ gives
\begin{equation}\label{eq:pmin18}
                         \mathcal P_{\min}(18)=26.
\end{equation}

\begin{lemma}[Equality reduction]\label{lem:equality}
If a $16$-tile $\{4,5\}$-polyform $A$ has two holes, then both holes consist
of one tile, $b(A)=15$, and $\overline A$ is an $18$-tile animal of
perimeter $26$.
\end{lemma}

\begin{proof}
Connectivity gives $b(A)\ge15$.  Filling does not change the outer
boundary, while each hole contains at least one tile.  The monotonicity used
in Lemma~2.1 of~\cite{EldridgeEtAl2026} (valid here because $p=4>3$) and
\eqref{eq:pmin18} give
\[
 p_o(A)=p_o(\overline A)
   \ge \mathcal P_{\min}(\lvert\overline A\rvert)
   \ge \mathcal P_{\min}(18)=26.
\]
Every one of the $64$ tile-edge incidences is either counted twice along a
shared edge or once on an outer or hole boundary.  Hence
\[
                  p_h(A)=64-2b(A)-p_o(A).
\]
Each of the two holes has boundary length at least four, so
\[
 8\le p_h(A)\le64-2\cdot15-26=8.
\]
All these inequalities are equalities.  Thus $b(A)=15$, $p_o(A)=26$, and
the two hole perimeters are both four.

The source convention regards the tiles in each hole as a polyform in their
own right.  A connected two-tile square animal has perimeter six, and
   $\mathcal P_{\min}$ is nondecreasing.  A hole containing at least two tiles would
therefore have perimeter at least six, contradicting its perimeter four.
Both holes are single tiles.  Exactly two tiles are added on filling, so
$\lvert\overline A\rvert=18$; its perimeter is $26$, which is minimal by
\eqref{eq:pmin18}.
\end{proof}

\section{The unique 18-tile extremal animal}

Malen, Rold\'an, and Toal\'a-Enr\'iquez define $C_{p,q}(1)$ to be an
edge-adjacent pair of tiles.  The next member $C_{p,q}(2)$ is obtained by
adding every tile incident to a perimeter vertex of this initial domino.
Their Theorem~6.4 states that, for every $k\ge1$, $C_{p,q}(k)$ is the unique
extremal animal at its tile count, up to isometry~\cite{MalenRoldanToala2023}.
The following local facts identify the relevant member of that subsequence.

\begin{lemma}[Domino-layer incidence]\label{lem:incidence}
The animal $C_{4,5}(2)$ has $18$ tiles, $23$ shared edges, and perimeter
$26$.  Its only tiles having all four edge-neighbours in $C_{4,5}(2)$ are
the two tiles of the initial domino, and those two tiles share an edge.
\end{lemma}

\begin{proof}
The domino has six boundary edges and six boundary vertices.  The six tiles
across its boundary edges are part of the new layer.  At each of the four
boundary vertices incident to one old tile, the cyclic star of five tiles
has two further positions after the old tile and the two across-edge tiles
are counted.  At each endpoint of the common edge, one further position
remains after the two old and two across-edge tiles are counted.  Thus the
new layer has
\[
          6+4(5-1-2)+2(5-2-2)=16
\]
tiles, for $18$ in total.

For completeness, all possible identifications and edge incidences in this
local count are checked by the exact ancillary verifier described below.
It shows that the $16$ new tiles form an edge-adjacency cycle.  Each of the
six across-edge tiles has three neighbours in $C_{4,5}(2)$, each of the ten
remaining new tiles has two, and each old domino tile has four.  Consequently
the sum of tile degrees is
\[
             2\cdot4+6\cdot3+10\cdot2=46,
\]
so there are $46/2=23$ shared edges and the perimeter is
$4\cdot18-2\cdot23=26$.  The displayed degrees also show that precisely the
old domino pair has no exposed edge, and the pair is edge-adjacent by its
definition.
\end{proof}

We now complete the mathematical argument.

\begin{proof}[Proof of Theorem~\ref{thm:main}]
Suppose that $A$ were a $16$-tile polyform with two holes.  By
Lemma~\ref{lem:equality}, $\overline A$ would be an $18$-tile extremal
animal of perimeter $26$.  Lemma~\ref{lem:incidence} shows that
$C_{4,5}(2)$ is an $18$-tile animal with that perimeter.  Theorem~6.4
of~\cite{MalenRoldanToala2023} therefore implies that $\overline A$ is
isometric to $C_{4,5}(2)$.

The two tiles in $\overline A\setminus A$ represent two distinct singleton
holes.  Each must have all four edge-neighbours still present; moreover,
the two deleted tiles cannot share an edge, since edge-adjacent absent tiles
belong to one complementary component.  They would thus have to be two
nonadjacent, fully edge-surrounded tiles of $C_{4,5}(2)$.  By
Lemma~\ref{lem:incidence}, the only two fully surrounded tiles are the
adjacent central domino.  Deleting them creates one two-tile cavity, not two
holes.  This contradiction proves the $16$-tile nonexistence.  The published
$17$-tile construction~\cite[Theorem~1.2 and Table~1]{EldridgeEtAl2026}
then gives $g_{4,5}(2)=17$.
\end{proof}

\section{Exact verification and scope}

The ancillary program \path{verify_c45_domino_layer.py} verifies only
the finite-incidence assertions in Lemma~\ref{lem:incidence}; the perimeter
reduction and the published uniqueness theorem remain human-readable parts
of the proof.  The program uses the standard faithful geometric
representation of the Coxeter group $[4,5]$~\cite[Chapter~5]{Humphreys1990}.
With simple reflections $s_0,s_1,s_2$, it checks
\[
 (s_0s_1)^4=(s_1s_2)^5=(s_0s_2)^2=1
\]
and represents every coordinate in
$\mathbb Q(\sqrt2,\sqrt5)$, with no floating-point comparisons.  It builds
a radius-four tile patch, selects exactly the tiles incident to the six
vertices of a central domino, and reconstructs all edge adjacencies.

\begin{center}
\begin{tabular}{@{}ll@{}}
\toprule
Checked quantity & Exact result \\
\midrule
tiles in $C_{4,5}(2)$ & $18$ \\
shared edges; perimeter & $23$; $26$ \\
new-layer edge graph & one $16$-cycle \\
fully edge-surrounded tiles & the two central tiles \\
central-pair relation & edge-adjacent \\
cell counts after deleting the pair $(V,E,F)$ & $(32,48,16)$ \\
bounded cavities after deleting the pair & $1$ \\
\bottomrule
\end{tabular}
\end{center}

Run the check from the ancillary directory with
\begin{verbatim}
python3 verify_c45_domino_layer.py
\end{verbatim}
under ordinary assertion-enabled Python.  The output serializes the complete
selected-tile adjacency lists and exact centre coordinates, so the stated
counts can be recomputed without geometric plotting.

This separation between an ambient theorem and the exact finite model that
a certificate actually checks is also explicit in the discrete-geometry
examples of Zeng et al.~\cite[Sections~6--7]{ZengEtAl2026Finite}.  In the
present proof, the certificate is not an exhaustive search over
$16$-tile polyforms.  Completeness instead comes from the equality reduction
and the published unique-extremal theorem; the computation only audits the
one forced animal.  Thus neither a successful program run nor a bounded
literature search is being used as a proof of the theorem's universal scope
or of publication priority.

\paragraph{Computational and AI assistance.}
OpenAI Codex assisted with primary-literature retrieval, formulation and
checking of the equality reduction, development and replay of the exact
finite-incidence verifier, and initial manuscript drafting.  Language-model
output is not mathematical evidence.  The proof rests on the cited source
theorems, the displayed deductions, and the deterministic exact verifier.
No proof assistant was used.

\paragraph{Data and code availability.}
The exact verifier is supplied with this article as
\path{verify_c45_domino_layer.py}.  It uses only the Python standard
library, takes no input, uses no randomness, and prints a JSON certificate.

\bibliographystyle{plainurl}
\bibliography{references}

\end{document}
